% answers.tex —— 1958 年·河北卷 · 参考答案与详细解析片段
% 由 tools/generate_projects.py 依来源总账自 1958/1958_河北_数学真题_试题.md + 1958/1958_河北_数学真题_解析正文.md + 1958/1958_河北_数学真题_答案要点.md 生成，勿手改（改源 MD 后重跑）


\gkpart{一、小题}

\begin{answer}
\gksolhead{1}{}

\textbf{题目}：

（本题满分 4 分）求多项式 $x^{376}-2x^{223}+5x^{145}-3x^{28}+3$ 用 $x+1$ 除所得的余数。

\Needspace{4\baselineskip}
\textbf{解析}

\Needspace{5\baselineskip}
多项式 $f(x)$ 除以 $x+1$ 的余数等于 $f(-1)$。代入：

\nopagebreak[4]
\[
f(-1)=(-1)^{376}-2(-1)^{223}+5(-1)^{145}-3(-1)^{28}+3
\]

\Needspace{5\baselineskip}
按指数的奇偶逐项计算：$(-1)^{376}=1$，$(-1)^{223}=-1$，$(-1)^{145}=-1$，$(-1)^{28}=1$，于是

\nopagebreak[4]
\[
f(-1)=1-2(-1)+5(-1)-3(1)+3
=1+2-5-3+3=-2
\]

故余数为 $-2$。

\gkanswer{答案：$-2$}

\end{answer}
\begin{answer}
\gksolhead{2}{}

\textbf{题目}：

（本题满分 6 分）试求函数 $y=\sin 2x$ 的周期，并作出它在一个周期内的图象。

\Needspace{4\baselineskip}
\textbf{解析}

\Needspace{5\baselineskip}
函数 $y=\sin 2x$ 的最小正周期为

\nopagebreak[4]
\[
T=\frac{2\pi}{2}=\pi
\]

在一个周期 $[0,\pi]$ 内，图象从 $(0,0)$ 出发，经 $\left(\dfrac\pi4,1\right)$、$\left(\dfrac\pi2,0\right)$、$\left(\dfrac{3\pi}{4},-1\right)$，回到 $(\pi,0)$，呈一条正弦曲线（频率为 $\sin x$ 的两倍）。

\gkanswer{答案：$\pi$}

\end{answer}
\begin{answer}
\gksolhead{3}{}

\textbf{题目}：

（本题满分 8 分）求 $\left(\dfrac{1}{\sqrt[3]{x}}+\sqrt{x}\right)^7$ 展开式中含 $x$ 的一次幂的项。

\Needspace{4\baselineskip}
\textbf{解析}

\Needspace{5\baselineskip}
把两项写成幂的形式：

\nopagebreak[4]
\[
\frac{1}{\sqrt[3]{x}}=x^{-1/3},\qquad \sqrt{x}=x^{1/2}
\]

\Needspace{5\baselineskip}
由二项式定理，通项为

\nopagebreak[4]
\[
C_7^k\left(x^{-1/3}\right)^{7-k}\left(x^{1/2}\right)^k
=C_7^k\,x^{-\frac{7-k}{3}+\frac{k}{2}}
\]

\Needspace{5\baselineskip}
令 $x$ 的指数等于 $1$：

\nopagebreak[4]
\[
-\frac{7-k}{3}+\frac{k}{2}=1
\]

\Needspace{5\baselineskip}
两边乘 $6$：$-2(7-k)+3k=6$，即 $-14+2k+3k=6$，$5k=20$，得 $k=4$。故含 $x$ 一次幂的项为

\nopagebreak[4]
\[
C_7^4x=35x
\]

\gkanswer{答案：$35x$}

\end{answer}
\begin{answer}
\gksolhead{4}{}

\textbf{题目}：

（本题满分 12 分）设一圆锥的母线长为 $l$，底面半径为 $r$，用平行于底面的一个平面在这圆锥上截下一个小圆锥，若小圆锥的侧面积等于原来圆锥的侧面积的一半，试求小圆锥的母线的长。

\Needspace{4\baselineskip}
\textbf{解析}

\Needspace{5\baselineskip}
设截下的小圆锥母线长为 $l'$。小圆锥与原圆锥相似，相似比为

\nopagebreak[4]
\[
\lambda=\frac{l'}{l}
\]

\Needspace{5\baselineskip}
于是小圆锥的底面半径为 $\lambda r$。小圆锥侧面积为

\nopagebreak[4]
\[
S_{\text{小}}=\pi(\lambda r)(\lambda l)=\lambda^2\pi rl
\]

\Needspace{5\baselineskip}
原圆锥侧面积为 $S=\pi rl$。由题意 $S_{\text{\text{小}}}=\dfrac12S$，得

\nopagebreak[4]
\[
\lambda^2\pi rl=\frac12\pi rl
\]

\Needspace{5\baselineskip}
故

\nopagebreak[4]
\[
\lambda^2=\frac12,\qquad \lambda=\frac{\sqrt2}{2}
\]

\Needspace{5\baselineskip}
因此小圆锥母线长

\nopagebreak[4]
\[
l'=\lambda l=\frac{\sqrt2}{2}l
\]

\gkanswer{答案：$\dfrac{\sqrt2}{2}l$}

\end{answer}
\gkpart{二、大题}

\begin{answer}
\gksolhead{5}{}

\textbf{题目}：

（本题满分 16 分）设 $O$ 为 $\triangle ABC$ 内的一点，线段 $AO$ 的延长线交 $BC$ 于 $A_1$ 点，线段 $BO$ 的延长线交 $AC$ 于 $B_1$ 点，线段 $CO$ 的延长线交 $AB$ 于 $C_1$ 点，求证：

（1）$\dfrac{\triangle OBC}{\triangle ABC}=\dfrac{OA_1}{AA_1}$；

（2）$\dfrac{OA_1}{AA_1}+\dfrac{OB_1}{BB_1}+\dfrac{OC_1}{CC_1}=1$。

\Needspace{4\baselineskip}
\textbf{解析}

\Needspace{5\baselineskip}
（1）$\triangle OBC$ 与 $\triangle ABC$ 有公共底边 $BC$，其面积之比等于它们到 $BC$ 的高之比。设 $A$ 到 $BC$ 的高为 $AA_1$ 方向上的垂线段长 $h_A$，$O$ 到 $BC$ 的高为 $h_O$。由于 $O$ 在线段 $AA_1$ 上，且 $A_1\in BC$，两个高之比等于 $OA_1:AA_1$，所以

\nopagebreak[4]
\[
\frac{\triangle OBC}{\triangle ABC}=\frac{h_O}{h_A}=\frac{OA_1}{AA_1}
\]

\Needspace{5\baselineskip}
（2）同理有

\nopagebreak[4]
\[
\frac{\triangle OCA}{\triangle ABC}=\frac{OB_1}{BB_1},\qquad
\frac{\triangle OAB}{\triangle ABC}=\frac{OC_1}{CC_1}
\]

\Needspace{5\baselineskip}
三式相加：

\nopagebreak[4]
\[
\frac{OA_1}{AA_1}+\frac{OB_1}{BB_1}+\frac{OC_1}{CC_1}
=\frac{\triangle OBC+\triangle OCA+\triangle OAB}{\triangle ABC}
\]

\Needspace{5\baselineskip}
而 $\triangle OBC$、$\triangle OCA$、$\triangle OAB$ 恰好把 $\triangle ABC$ 无重叠、无遗漏地分成三部分，故其面积和等于 $\triangle ABC$。于是

\nopagebreak[4]
\[
\frac{OA_1}{AA_1}+\frac{OB_1}{BB_1}+\frac{OC_1}{CC_1}=1
\]

\gkanswer{答案：见解析版}

\end{answer}
\begin{answer}
\gksolhead{6}{}

\textbf{题目}：

\Needspace{5\baselineskip}
（本题满分 18 分）解方程组：

\nopagebreak[4]
\[
\begin{cases}
\cos 2x+\cos 2y=\dfrac{\sqrt3}{2},\\[2pt]
x+y=\arccos\dfrac12.
\end{cases}
\]

\Needspace{4\baselineskip}
\textbf{解析}

\Needspace{5\baselineskip}
由 $x+y=\arccos\dfrac12=\dfrac\pi3$，得 $x+y=\dfrac\pi3$。利用和差化积：

\nopagebreak[4]
\[
\cos2x+\cos2y=2\cos(x+y)\cos(x-y)
=2\cos\frac\pi3\cos(x-y)=\cos(x-y)
\]

\Needspace{5\baselineskip}
代入原方程得

\nopagebreak[4]
\[
\cos(x-y)=\frac{\sqrt3}{2}
\]

\Needspace{5\baselineskip}
故

\nopagebreak[4]
\[
x-y=\pm\frac\pi6+2k\pi\qquad(k\in\mathbb Z)
\]

与 $x+y=\dfrac\pi3$ 联立：

\begin{itemize}
\item 当 $x-y=\dfrac\pi6+2k\pi$ 时，$2x=\dfrac\pi2+2k\pi$，得 $x=\dfrac\pi4+k\pi$，$y=\dfrac\pi{12}-k\pi$；
\item 当 $x-y=-\dfrac\pi6+2k\pi$ 时，$2x=\dfrac\pi6+2k\pi$，得 $x=\dfrac\pi{12}+k\pi$，$y=\dfrac\pi4-k\pi$。
\end{itemize}

\gkanswer{答案：$\begin{cases}x=\dfrac\pi4+k\pi,\\ y=\dfrac\pi{12}-k\pi\end{cases}$ 或 $\begin{cases}x=\dfrac\pi{12}+k\pi,\\ y=\dfrac\pi4-k\pi\end{cases}$（$k$ 为整数）}

\end{answer}
\begin{answer}
\gksolhead{7}{}

\textbf{题目}：

（本题满分 18 分）一农场有甲、乙两台打麦机，甲机的工作效率是乙机的 $2$ 倍，先只用甲机打完全部麦子的 $\dfrac23$，然后只用乙机继续打完，所需要的时间比同时用两台打麦机打完全部麦子所需要的时间多 $4$ 天，问分别用一台打麦机打完所有的麦子各需要多少天？

\Needspace{4\baselineskip}
\textbf{解析}

设乙机单独打完所有麦子需要 $x$ 天，则乙机每天打 $\dfrac1x$。甲机效率是乙机的 $2$ 倍，故甲机单独打完需要 $\dfrac{x}{2}$ 天，每天打 $\dfrac2x$。

\Needspace{5\baselineskip}
两机同时打，每天共打 $\dfrac2x+\dfrac1x=\dfrac3x$，同时打完需要

\nopagebreak[4]
\[
\frac{x}{3}\ \text{天}
\]

\Needspace{5\baselineskip}
先只用甲机打完全部麦子的 $\dfrac23$，需要

\nopagebreak[4]
\[
\frac{2/3}{2/x}=\frac{x}{3}\ \text{天}
\]

\Needspace{5\baselineskip}
再用乙机打完剩余的 $\dfrac13$，需要

\nopagebreak[4]
\[
\frac{1/3}{1/x}=\frac{x}{3}\ \text{天}
\]

\Needspace{5\baselineskip}
所以先甲后乙共需 $\dfrac{2x}{3}$ 天。由题意比同时打多 $4$ 天：

\nopagebreak[4]
\[
\frac{2x}{3}=\frac{x}{3}+4
\]

解得 $x=12$。故乙机单独需 $12$ 天，甲机单独需 $\dfrac{12}{2}=6$ 天。

\gkanswer{答案：甲机 $6$ 天，乙机 $12$ 天}

\end{answer}
\begin{answer}
\gksolhead{8}{}

\textbf{题目}：

（本题满分 18 分）已知方程 $4x^2-(4\cos\theta)x+1=0$ 有实根，求 $\theta$ 的一切可能的值。

\Needspace{4\baselineskip}
\textbf{解析}

\Needspace{5\baselineskip}
方程 $4x^2-(4\cos\theta)x+1=0$ 有实根的充要条件是判别式不小于零：

\nopagebreak[4]
\[
\Delta=(4\cos\theta)^2-4\cdot4\cdot1=16\cos^2\theta-16=16(\cos^2\theta-1)\ge0
\]

\Needspace{5\baselineskip}
故

\nopagebreak[4]
\[
\cos^2\theta-1\ge0
\]

\Needspace{5\baselineskip}
而 $\cos^2\theta\le1$ 恒成立，故只能 $\cos^2\theta=1$，即 $\cos\theta=\pm1$。因此

\nopagebreak[4]
\[
\theta=n\pi\qquad(n\in\mathbb Z)
\]

\gkanswer{答案：$\theta=n\pi$（$n$ 为整数）}

\end{answer}